\section{Podsumowanie}

% SLIDE ID: sec11-goal
\BlockGoalFrame{\SlideID{sec11-goal}}{
  \item zebrać najważniejsze wnioski z całego dnia,
  \item jeszcze raz połączyć systemy z logiką procesu badania,
  \item domknąć wykład pytaniami i dyskusją końcową.
}

% SLIDE ID: sec11-key-takeaways
\begin{frame}{5 najważniejszych wniosków z całego dnia}
  \SlideID{sec11-key-takeaways}
  \begin{itemize}
    \item Badanie kliniczne to jeden proces wspierany przez wiele wyspecjalizowanych systemów.
    \item \texttt{EDC} jest centrum danych klinicznych, ale nie zastępuje \texttt{CTMS}, \texttt{IWRS} ani \texttt{eTMF}.
    \item Jakość danych, bezpieczeństwo i gotowość inspekcyjna powstają na styku ludzi, procesów i technologii.
    \item Dane od pacjenta, integracje i \texttt{AI} zwiększają możliwości, ale też wymagają większej kontroli.
    \item Najlepsze systemy nie usuwają odpowiedzialności człowieka, tylko pomagają ją lepiej realizować.
  \end{itemize}
\end{frame}

% SLIDE ID: sec11-system-map
\begin{frame}{Mapa systemów jeszcze raz}
  \SlideID{sec11-system-map}
  \begin{itemize}
    \item \texttt{CTMS}: zarządzanie badaniem, ośrodkami, monitoringiem i oversight.
    \item \texttt{eTMF}: dokumentacja, wersje, gotowość audytowa i inspekcyjna.
    \item \texttt{EDC / eCRF}: dane kliniczne, walidacje, \texttt{query}, cleaning i \texttt{database lock}.
    \item \texttt{IWRS / IVRS}: randomizacja i lek badany.
    \item \texttt{ePRO}: dane bezpośrednio od pacjenta.
    \item Safety + \texttt{AI}: bezpieczeństwo, sygnały ryzyka i wsparcie analityczne.
  \end{itemize}
\end{frame}

% SLIDE ID: sec11-learning-outcomes
\begin{frame}{Co uczestnik powinien umieć po zajęciach}
  \SlideID{sec11-learning-outcomes}
  \begin{itemize}
    \item odróżnić role głównych systemów \texttt{eClinical},
    \item opisać podstawowy przepływ danych od pacjenta do \texttt{database lock},
    \item rozpoznać miejsca krytyczne dla jakości danych i bezpieczeństwa,
    \item rozumieć, gdzie kończy się automatyzacja, a zaczyna odpowiedzialność człowieka,
    \item lepiej współpracować z zespołami operacyjnymi, \texttt{DM}, monitoringiem i safety.
  \end{itemize}
\end{frame}

% SLIDE ID: sec11-control-questions-1
\begin{frame}{Pytania kontrolne 1}
  \SlideID{sec11-control-questions-1}
  \begin{enumerate}
    \item Czym różni się \texttt{EDC} od \texttt{CTMS}?
    \item Po co w \texttt{EDC} \texttt{audit trail}?
    \item Co jest największym ryzykiem w \texttt{IWRS}?
  \end{enumerate}
\end{frame}

% SLIDE ID: sec11-control-questions-2
\begin{frame}{Pytania kontrolne 2}
  \SlideID{sec11-control-questions-2}
  \begin{enumerate}
    \setcounter{enumi}{3}
    \item Dlaczego \texttt{eTMF} nie jest tylko repozytorium plików?
    \item Jak \texttt{ePRO} zmienia rolę pacjenta?
    \item Gdzie \texttt{AI} realnie pomaga w badaniu klinicznym?
  \end{enumerate}
\end{frame}

% SLIDE ID: sec11-discussion
\DiscussionFrame{Dyskusja końcowa}{
  Który system wydaje się dziś najbardziej krytyczny z perspektywy jakości badania? W którym miejscu procesu najłatwiej o błąd i co jako zespół można zrobić, żeby to ryzyko ograniczyć?
}

% SLIDE ID: sec11-contact
\begin{frame}{Dziękuję}
  \SlideID{sec11-contact}
  \begin{center}
    {\Large Nowe technologie i innowacje w badaniach klinicznych}\par
    \vspace{1em}
    {\large dr Krzysztrof Misztal}\par
    \vspace{1.5em}
    \begin{beamercolorbox}[wd=0.72\linewidth,sep=1em,rounded=true]{block body}
      \centering
      \textbf{Miejsce na kontakt}\par
      \vspace{0.5em}
      e-mail / LinkedIn / dane uczelni
    \end{beamercolorbox}
  \end{center}
\end{frame}
